\section{MLP-Mixer: un experimento de control que debilita aún más el sesgo inductivo}

Al final, la clase extiende con MLP-Mixer la idea de «dejar que los datos hablen». Mixer no usa convolution ni self-attention: aplica MLP de forma alternada sobre la dimensión de token y la dimensión de channel. No pretende demostrar que attention sea inútil, sino comprobar si, con preentrenamiento a gran escala, un dense mixing simple puede aprender representaciones visuales suficientemente buenas.

\subsection{De ViT a Mixer: una pregunta de investigación, no un cambio de marca}

La diapositiva de transición reinterpreta el éxito de ViT como «quizá la escala y los datos importen más que un operator específico». Por eso Mixer es una controlled provocation: conserva el patch embedding, la residual/normalization y la recipe a gran escala, sustituye el token interaction mechanism y observa si la transfer scaling se mantiene.

\lecturefigure{slide-34-mixer-divider.jpg}{MLP-Mixer continúa con «letting the data speak» y retira además attention.}{Video oficial, 1:03:36.}

\begin{knowledgebox}{Lectura de la figura: no concluir que «la arquitectura no importa»}
Si Mixer resulta competitivo con grandes volúmenes de datos, eso solo indica que attention no es el único mecanismo para obtener representaciones sólidas. La estabilidad del entrenamiento, la compute efficiency, la resolution scaling, la dense prediction y el desempeño con pocos datos todavía pueden diferir. La conclusión más razonable es que existen varias combinaciones viables de operator, scale y recipe, y que hay que elegir según la tarea y las restricciones del sistema.
\end{knowledgebox}

\subsection{Token-mixing y channel-mixing}

En la matriz de entrada $X\in\mathbb{R}^{N\times D}$, $N$ son los patch tokens y $D$ son los channels. La token-mixing MLP mezcla, para cada channel, a través de las $N$ posiciones; la channel-mixing MLP mezcla, para cada token, a través de los $D$ canales. Ambas se alternan junto con residual connection, de modo que la información puede propagarse entre posiciones espaciales y también combinarse en la dimensión de características.

\lecturefigure{slide-35-mlp-mixer-architecture.jpg}{MLP-Mixer alterna las MLP de token-mixing y channel-mixing.}{Video oficial, 1:03:44; Tolstikhin et al. 2021.}

\begin{importantbox}{Forma abstracta del bloque Mixer}
Si se ignoran los detalles de LayerNorm, puede escribirse como
\[
U=X+W_2\,\sigma(W_1X^{\top})^{\top},\qquad
Y=U+\sigma(UV_1)V_2.
\]
La primera expresión mezcla a lo largo de la dimensión de token y la segunda a lo largo de la dimensión de channel; $\sigma$ es una no linealidad. A diferencia de attention, los pesos de token mixing no dependen del contenido de la entrada actual, por lo que la interacción es estática; el costo y las propiedades de generalización también varían con el número de tokens y el tamaño de las matrices de pesos.
\end{importantbox}

\subsection{Scaling comparison: la pendiente de una familia de modelos importa más que un punto aislado}

La última figura compara el desempeño linear few-shot de Mixer, ViT y BiT/ResNet a medida que crecen los datos de entrenamiento. Lo importante es que las tres clases de modelos se benefician de la escala, pero difieren en la pendiente, en el punto de partida con pocos datos y en el techo con muchos datos; esto respalda de nuevo el marco según el cual «el prior arquitectónico determina la sample efficiency y la escala de los datos determina si un modelo flexible logra alcanzar a los demás».

\lecturefigure{slide-36-mixer-scaling.jpg}{Linear few-shot scaling de Mixer, ViT y BiT con distintos tamaños del conjunto de entrenamiento.}{Video oficial, 1:05:32.}

\begin{knowledgebox}{Lectura de la figura: comparar el punto de partida, la pendiente y la zona de saturación}
En el extremo de pocos datos, los modelos con un prior visual fuerte suelen partir de un punto más alto; conforme aumentan los datos de entrenamiento, las curvas de ViT/Mixer pueden ser más pronunciadas; que sigan a la cabeza en el extremo de muchos datos depende de la capacidad y de la recipe. Un solo punto con la mayor cantidad de datos no demuestra sample efficiency, y un solo punto con la menor cantidad tampoco demuestra asymptotic capability. Solo la curva completa muestra la relación de intercambio entre inductive bias y scale.
\end{knowledgebox}

\teachervoice{El ponente coloca Mixer al final no para anunciar que «MLP reemplaza a Transformer», sino para subrayar que la investigación debe seguir desglosando qué contribuyó realmente al resultado: attention, la patch representation, la escala, los datos o el régimen de entrenamiento. La postura que el docente expresa de viva voz es más cautelosa que el título de la slide.}

\subsection{Resumen de la sección}

MLP-Mixer demuestra que attention no es el único candidato para obtener representaciones visuales a gran escala. Lleva el hilo conductor de la clase a una conclusión más general: entre el sesgo inductivo, la escala de los datos y la eficiencia del sistema existe un intercambio, y las familias de modelos deben compararse mediante la scaling curve completa y la transfer en múltiples tareas.
